\chapter{Condiciones de frontera de las ecuaciones de Maxwell en dieléctricos}
\label{sec:boundary_conditions_in_dielectrics}


\begin{figure}[h!]
    \centering
    % Placeholder for the image of the diagram. The actual image would be inserted here.
\end{figure}

Partamos de las ecuaciones de Maxwell para un dieléctrico caracterizado por 
las constantes dieléctricas $(\epsilon, \mu)$ 

\begin{equation}
    \begin{gathered}
        \div \mathbf{E} = \frac{\rho}{\epsilon} \\
        \curl \mathbf{E} = -\pdv{\mathbf{B}}{t}
    \end{gathered}
    \qquad
    \begin{gathered}
        \div \mathbf{B} = 0 \\
        \curl \mathbf{B} = \mu \mathbf{J} + \frac{1}{v^2}\pdv{\mathbf{E}}{t}
    \end{gathered}
    \qquad
    v = \sqrt{\frac{1}{\mu\epsilon}}
    \label{eq:maxwell_dielectric_basic}
\end{equation}

O más apropiadamente para este análisis en términos de los campos auxiliares $\mathbf D$ y $\mathbf H$ definidos como 

\begin{align}
\mathbf{D} &= \varepsilon \mathbf{E}, \\[6pt]
\mathbf{B} &= \mu \mathbf{H},
\end{align}

las ecuaciones de Maxwell quedan como 

\begin{align}
\nabla \cdot \mathbf{D} &= \rho_f, 
&\qquad \nabla \cdot \mathbf{B} &= 0, \\[6pt]
\nabla \times \mathbf{E} &= -\,\frac{\partial \mathbf{B}}{\partial t}, 
&\qquad \nabla \times \mathbf{H} &= \mathbf{J}_f + \frac{\partial \mathbf{D}}{\partial t}.
\label{eq:maxwell_aux_dielectric}
\end{align}

Sea un medio $M$ y otro $M'$ separados por una interfaz $\Sigma: \{ (x,y,z) | z=0 \}$, donde $M$ es el medio en $z<0$ y el medio $M'$ es el medio en $z>0$. Estamos interesados en relacionar el campo $\mathbf{E}$ y $\mathbf{H}$ en el medio $M'$ a partir de los campos $\mathbf{E}$ y $\mathbf{B}$ en el medio $M$. 

Para facilitar nuestro análisis, definimos este campo arbitrario $\mathbf{F}(\mathbf r,t)$ de la forma:

\begin{equation}
    \mathbf{F}(\mathbf r,t) = \mathbf{F}(\mathbf r,t)\Theta(-z) + \mathbf{F}'(\mathbf r,t)\Theta(z)
    \label{eq:field_piecewise}
\end{equation}

donde

\begin{equation}
    \Theta(z) =
    \begin{cases}
        1 & z>0 \\
        0 & \text{altura}
    \end{cases}, 
\end{equation}

de donde se sigue que el campo sin prima representa el campo en el medio $M$ y el primado el campo en el medio $M'$. 

\text{Calculamos ahora el $\div$ y el $\curl$ del campo $\mathbf{F}$.}

\begin{equation} 
    \div \mathbf{F}(\mathbf r,t) = \nabla \cdot \mathbf{F}
\end{equation}

\begin{equation}
    \begin{split}
        &= [\nabla \cdot \mathbf{F}(\mathbf r,t)] \Theta(-z) + \mathbf{F}(\mathbf r,t) \cdot \nabla \Theta(-z) \\
        &+ [\nabla \cdot \mathbf{F}'(\mathbf r,t)] \Theta(z) + \mathbf{F}'(\mathbf r,t) \cdot \nabla \Theta(z)
    \end{split}
\end{equation}

\begin{equation}
    \nabla \cdot \mathbf{F} = [\nabla \cdot \mathbf{F}] + \delta(z) [\mathbf{F}'_z - \mathbf{F}_z]
    \label{eq:div_field}
\end{equation}

\pagebreak

Para el rotacional:

\begin{equation}
    \nabla \times \mathbf{F} = [\nabla \times \mathbf{F}] + \delta(z) \hat{n} \times (\mathbf{F}' - \mathbf{F})
    \label{eq:curl_field}
\end{equation}

Escribimos ahora las fuentes en una descomposición para la superficie
y para fuera de esta como

\begin{equation}
    \rho = \rho_f + \sigma(x,y,t)\delta(z), 
    \qquad 
    \mathbf{J} = \mathbf{J}_f + \mathbf{K}(x,y,t)\delta(z)
    \label{eq:surface_sources}
\end{equation}

Escribimos las ecuaciones de Maxwell en la forma:

\begin{equation}
    \begin{gathered}
        \div \mathbf{D} = \rho_f \\
        \curl \mathbf{E} = -\pdv{\mathbf{B}}{t}
    \end{gathered}
    \qquad
    \begin{gathered}
        \div \mathbf{B} = 0 \\
        \curl \mathbf{H} = \mathbf{J}_f + \pdv{\mathbf{D}}{t}
    \end{gathered}
    \label{eq:maxwell_form_expanded}
\end{equation}

Podemos entonces, por ejemplo, usar las expresiones \eqref{eq:div_field}, \eqref{eq:curl_field} y \eqref{eq:surface_sources} en las ecuaciones \eqref{eq:maxwell_form_expanded}.  
Al cancelar los términos proporcionales a $\delta(z)$, se obtienen las condiciones de frontera para que se cumplan las ecuaciones de Maxwell. Con lo que se llega a:

\begin{equation}
    \begin{gathered}
        \hat{n} \cdot (\mathbf{D}_2 - \mathbf{D}_1) = \sigma_f \\
        \hat{n} \cdot (\mathbf{B}_2 - \mathbf{B}_1) = 0 \\
        \hat{n} \times (\mathbf{E}_2 - \mathbf{E}_1) = 0 \\
        \hat{n} \times (\mathbf{H}_2 - \mathbf{H}_1) = \mathbf{K}_f
    \end{gathered}
    \label{eq:boundary_conditions}
\end{equation}
